% Generated draft: requires mathematical and source review.
\documentclass[11pt]{article}
\usepackage[margin=25mm]{geometry}
\usepackage{fontspec}
\setmainfont{Latin Modern Roman}
\usepackage{amsmath,amssymb,mathtools,unicode-math}
\setmathfont{Latin Modern Math}
\usepackage{enumitem,needspace,xurl,hyperref}
\setlength{\parindent}{0pt}
\setlength{\parskip}{5pt}
\setlength{\emergencystretch}{3em}
\newcommand{\N}{\mathbb{N}}
\newcommand{\Z}{\mathbb{Z}}
\newcommand{\Q}{\mathbb{Q}}
\newcommand{\R}{\mathbb{R}}
\newcommand{\C}{\mathbb{C}}
\newcommand{\F}{\mathbb{F}}
\newcommand{\A}{\mathbb{A}}
\newcommand{\PP}{\mathbb{P}}
\newcommand{\E}{\mathbb{E}}
\newcommand{\eps}{\varepsilon}
\newcommand{\dd}{\,\mathrm{d}}
\begin{document}
\section*{UnsolvedMath ID: 30004832}
\textbf{Source problem code:} catalogue entry
\subsection{Definitions and mathematical statement}
For terminology only, a standard setting for KMS states consists of a unital \(C^*\)-algebra \(A\), a one-parameter group \(\alpha:\R\to\operatorname{Aut}(A)\) such that \(t\mapsto\alpha_t(a)\) is norm-continuous for every \(a\in A\), and an inverse temperature \(\beta>0\). A state is a positive linear functional \(\varphi:A\to\C\) with \(\varphi(1)=1\). Such a state satisfies the \(\beta\)-KMS condition if, for every \(a,b\in A\), there is a bounded continuous function \(F_{a,b}\) on \(\{z\in\C:0\leq\operatorname{Im}z\leq\beta\}\), holomorphic in its interior, such that
\[
F_{a,b}(t)=\varphi(a\alpha_t(b)),\qquad
F_{a,b}(t+i\beta)=\varphi(\alpha_t(b)a)
\quad\text{for every }t\in\R.
\]
This definition supplies terminology only: the extracted source does not specify this setting, an algebra, a dynamics, an inverse temperature, or any categorical input.

The supplied question is: can further information about the structure of KMS states be derived from purely categorical considerations?

Here both the categorical data and the additional structure to be determined remain unspecified. In particular, the excerpt does not identify a category, any extra categorical structure, its relationship to a dynamical system, the information already known, or the intended quantifiers. Consequently, the supplied evidence supports this research question but does not determine a precise universally or existentially quantified mathematical assertion. Choosing such data or a target invariant would require additional source context.
\subsection{Short English statement}
The question asks whether categorical information can reveal additional structure of KMS states. The supplied excerpt does not identify the relevant category, dynamical system, temperature, or desired structural conclusion, so a precise theorem or conjecture cannot yet be extracted. Its open status is supplied dated curation and has not been independently verified.
\subsection{Sources}
\begin{itemize}
\item Local UnsolvedMath registry, supplied complete source record: permanent ID \(30004832\), title \emph{Categorical Determination of KMS-State Structure}, exact source problem code \texttt{catalogue entry}, in the Oberwolfach Reports Open Problems collection. No individual proposer is supplied.
\item Supplied source citation: \emph{Quantum Groups -- Algebra, Analysis and Category Theory} (2022), Oberwolfach Reports, DOI \texttt{10.4171/owr/2021/44}, pp.\ 2397--2458. The extracted text carries the page marker 2419. \url{https://doi.org/10.4171/owr/2021/44}. The registry assigns proposed year 2021; the relationship between that year and the citation's 2022 date is not explained.
\item Supplied literature source: \emph{Categorical compact quantum groups}. \url{https://arxiv.org/abs/1603.05519}. The record does not specify a passage establishing the connection to this KMS-state question.
\item Literature assessment dated exactly \texttt{2026-08-22}: the supplied curation labels the problem open and states that the extracted sentence leaves open whether categorical data alone determine further structure of KMS states, and that current categorical classification results do not provide such a general implication. This is supplied dated curation, not an independently verified literature assessment. The background explicitly requests independent verification of time-sensitive status and cited claims.
\end{itemize}
\end{document}
